% !TeX root = ../main.tex

\chapter{Crystal Symmetry \& Crystal Structure}

\section{Crystal Symmetry}

\subsection{Inversion Symmetry}

\subsection{Rotoinversion Symmetry}

\section{Crystal Structure}

\begin{equation}
  \text{Crystal lattice} + \text{Basis} = \text{Crystal structure}
\end{equation}
Replace a every site (represent the symmetry of crystals) by basis (represent the compound of crystal: one / a group of atom / a molecule / a set of molecule): nvotif

Fourier transfor:
\begin{equation}
  \upe^{\iu\bm G\cdot\bm R} = 1,~
  A(k) = A(k) \upe^{\iu\bm k\cdot\bm R}
\end{equation}
then $\upe^{\iu\bm k\cdot\bm R} = 1$.

\begin{equation}
  \bm a_i \cdot \bm b_i = 2\pi\delta_{ij}
\end{equation}
\begin{example}[Rec. lattice of BCC]
  
\end{example}

\subsection{Structure factor of a monatomic lattice}

In BCC structure, there are two atoms in a unit cell: $(0,0,0)$ and $\frac a2(1,1,1)$.
They have the same atomic form factor $f$
\[
  S_G = f \upe^{-\iu \bm G_1 \cdot \bm r_1} +
        f \upe^{-\iu \bm G_2 \cdot \bm r_2}
\]
where $\bm G_1 = \frac{2\pi}{a} (hx_0 + kx_0 + lx_0)$ and $\bm G_1 \cdot \bm r_1 = 0$;
$\bm G_2 = \frac{2\pi}{a} (1 \times \frac a2h + 1 \times \frac a2k + 1 \times \frac a2l)$.
\begin{itemize}
  \item $S_G = 0$, when $h + k + l$ is an odd integer, there No diffraction.
  \item $S_G = 2f$, when $h + k + l$ is an even integer, diffraction peak exists.
\end{itemize}

X-Ray Diffraction: some peeks in $I$ along $\theta$: stands for $(100)$, $(111)$, $(110)$, but $(110)$ is missing.
\emph{Missing orders make it easier to identify a particular crystal structure.}

The wave length of the X-Ray source has a range $(\lambda_1, \lambda_2)$,
corresponging $k_{\max}$ to $k_{\min}$.
For each $k$, we have a circle `field'.

For incident beam, we can get a series of ring: the distance between the rings
$d = a/\sqrt{h^2 + k^2 + l^2}$ for cubic lattice.

\subsection{Bloch theorem}

Lattice vector
\[
  \bm R = k_1 \bm x_1 + k_2 \bm x_2 + k_3 \bm x_3
\]
where $k_i$ should be integer numbers.
Bloch electrons
\[
  \ab[-\frac{\hbar^2}{2m} \nabla^2 + U(\bm r)] \psi(\bm r) = E\psi(\bm r)
\]
where $U(\bm r + \bm R) = u(\bm r)$.
The solution should be the form of
\begin{equation}
  \psi_k(\bm r) = \upe^{\iu\bm k \cdot \bm r} u_k(\bm r),
\end{equation}
where $u_k(\bm r) = u_k(\bm r + \bm R)$,
multiplication of Plane wave and lattice periodic function.

\subsection{Nearly free electron model}

For free electrons, the periodic function
\begin{equation}
  E = \frac{\hbar^2k^2}{2m}
\end{equation}
then, when we apply the Bloch theorem
\[
  \psi_R(\bm r) = \upe^{\iu \bm k \cdot \bm r} u_k(\bm r), \quad
  u_k(\bm r) = u_k(\bm r + \bm R)
\]

\begin{example}[Simple cube]
  \[
    E(\bm k) = \epsilon_s - J_0 - 2J_1(\cos(k_xa) + \cos(k_ya) + \cos(k_za))
  \]
  where
  \[
    -\bm k \cdot \bm R = -k_x \cdot a + ky \cdot 0 + k_z \cdot 0 = -k_x \cdot a
  \]
  In the first Brillouin zone of simple cube, there are several high symmetric
  points
  \begin{enumext}[columns = 2]
    \item $\Gamma$ points: $\bm k = (0,0,0)$
    \item $X$ points: $\bm k = (\pi/a,0,0)$
    \item $R$ points: $\bm k = (\pi/a,\pi/a,\pi/a)$
    \item $M$ points: $\bm k = (\pi/a,\pi/a,0)$
  \end{enumext}
  Along different directions, the band structure is different.
  \begin{align}
    E(\Gamma) & = \epsilon_0 - J_0 - 6J_1,\\
    E(X)      & = \epsilon_0 - J_0 - 2J_1,\\
    E(R)      & = \epsilon_0 - J_0 + 6J_1,
  \end{align}
  The effective mass. Around the ``minimum'' point, the electrons behave as the
  free electrons, i.e.,
  \[
    E = \frac{\hbar^2k^2}{2m}
  \]
  Since we have (for cubic)
  \[
    E(k) = \epsilon_0 - J_0 - 6J_1 + J_1|\bm k|^2 a^2
  \]
  and we can expnad it into Taylor series
  \[
    E = E_0 + a_1\pdv Ek + a_2\pdv[2]Ek + \cdots
  \]
  where $\pdv Ek$ around the minimum point. Compare the expressions, we have the effective mass
  \[
    m^* = \frac{\hbar^2}{\pdv[2]E/k}
  \]
\end{example}
The bands have antisymmetric properties, i.e.,
\[
  \pdv E{k_x} \neq \pdv E{k_y} \neq \pdv E{k_z}, \qq{and}
  \pdv[2] E{k_x} \neq \pdv[2] E{k_y} \neq \pdv[2] E{k_z},
\]

\section{Bloch Oscillations}

For three dimensional
\[
  E = \epsilon_0 - J_0 - 2J_1(\cos(k_xa) + \cos(k_ya) + \cos(k_za))
\]
For one dimensional
\[
  E = \epsilon_0 - J_0 - 2J_1 \cos(k_xa)
\]

\subsection{Requirements for observation}

\[
  w = \frac{eEa}{\hbar}, \quad T = \frac{2\pi}{eEa} = \frac{h}{eEa}
\]